%%% ============================================================
%%%  package_and_settings.tex
%%%  統一管理所有套件載入與全域設定
%%%  【注意】請勿任意更動本檔，修改前請先備份
%%% ============================================================


% ─────────────────────────────────────────
%  一、版面設定
% ─────────────────────────────────────────
% fancyhdr：自訂頁首（header）與頁尾（footer）的內容與格式
\usepackage{fancyhdr}
% geometry：設定版面尺寸與各邊邊距
\usepackage{geometry}
\geometry{headheight=2.54cm, bottom=2.54cm, right=3.54cm, left=3.54cm}
% appendix：支援附錄（Appendix）環境，可自動調整附錄標題與編號
\usepackage{appendix}
% framed：提供 framed 環境，可在段落四周繪製簡單框線
\usepackage{framed}


% ─────────────────────────────────────────
%  二、圖片設定
% ─────────────────────────────────────────
% graphicx：插入外部圖片（\includegraphics），支援縮放、旋轉等
\usepackage{graphicx}
% float：控制浮動體（圖、表）排版位置，支援 [H] 強制原地定位
\usepackage{float}
% subfigure：在單一 figure 環境中並排多張子圖（可各自編號與加標題）
\usepackage{subfigure}
% 圖片路徑：本模板不設全域 graphicspath，
% 請在 \includegraphics 中使用從根目錄出發的完整相對路徑，例如：
%   \includegraphics[width=0.8\textwidth]
%     {content/__SUBJECT__/note/note_1/figures/my_diagram}
% 這樣無論在哪個文件資料夾，路徑都不會出錯。


% ─────────────────────────────────────────
%  三、表格設定
% ─────────────────────────────────────────
% booktabs：提供專業表格用的水平線指令（\toprule、\midrule、\bottomrule）
\usepackage{booktabs}
% longtable：支援跨頁的長表格
\usepackage{longtable}
% dcolumn：讓表格欄位以小數點（或指定字元）對齊
\usepackage{dcolumn}
% multirow：支援表格儲存格跨列合併（\multirow{}）
\usepackage{multirow}


% ─────────────────────────────────────────
%  四、TikZ / 繪圖設定
% ─────────────────────────────────────────
% tikz：向量繪圖核心套件，用於繪製幾何圖形、流程圖等
\usepackage{tikz}
\tikzset{>=latex}                      % 將所有箭頭預設為 LaTeX 樣式
% pgfplots：繪製函數圖形與資料圖，以 tikz 為基礎
\usepackage{pgfplots}
% xcolor：定義與使用顏色，供 tikz 及其他環境共用
\usepackage{xcolor}
% contour：在文字四周產生暈光（halo）效果，常用於 tikz 圖形標籤
\usepackage[outline]{contour}
\contourlength{1.2pt}
% 常用 tikz 函式庫
\usetikzlibrary{positioning, calc}    % 相對定位與坐標計算
\usetikzlibrary{backgrounds}          % 繪製背景圖層（inner frame sep 需要）
\usetikzlibrary{arrows.meta, calc}    % 進階箭頭樣式
\usetikzlibrary{patterns}             % 填充圖樣（斜線、點陣等）
% pgfplots 函式庫：填充兩函數之間的區域（積分示意圖）
\usepgfplotslibrary{fillbetween}
\pgfplotsset{compat=1.12}             % 確保 TikZ 坐標與軸坐標互換的相容性

% mathtools：用於公式排版，需在 tkz-euclide 之前載入
\usepackage{mathtools}
% tkz-euclide：歐幾里得幾何繪圖工具（標記點、角度等）
\usepackage{tkz-euclide}

% 自訂 tkzPointShowCoord 的箭頭與座標標籤樣式
\makeatletter
\tikzset{arrow coord style/.style={%
        densely dashed,
        \tkz@euc@linecolor,
}}
\tikzset{xcoord style/.style={%
        \tkz@euc@labelcolor,
        font=\normalsize, text height=1ex,
        inner sep=0pt, outer sep=0pt,
        fill=\tkz@fillcolor,
        below=6pt
}}
\tikzset{ycoord style/.style={%
        \tkz@euc@labelcolor,
        font=\normalsize, text height=1ex,
        inner sep=0pt, outer sep=0pt,
        fill=\tkz@fillcolor,
        left=6pt
}}
\makeatother


% ─────────────────────────────────────────
%  五、程式碼嵌入設定
% ─────────────────────────────────────────
% listings：在文件中插入原始碼區塊（支援語法高亮）
% lstautogobble：自動去除 lstlisting 環境中的縮排空白
\usepackage{listings, lstautogobble}
\lstset{
    backgroundcolor = \color{lightgray},    % 背景色（依賴第十區段定義的 lightgray）
    basicstyle      = {\small\ttfamily},    % 預設字型（等寬小字）
    breaklines      = true,                 % 超出行寬時自動換行
    postbreak       = \raisebox{0ex}[0ex][0ex]{\ensuremath{\color{lightgray}\hookrightarrow\space}},
    stringstyle     = \color{darkcyan},     % 字串顏色（依賴 darkcyan）
    otherkeywords   = {0,1,2,3,4,5,6,7,8,9},
    morekeywords    = {TRUE, FALSE},
    deletekeywords  = {data, frame, length, as, character},
    keywordstyle    = \color{blue},         % 關鍵字顏色
    commentstyle    = \color{darkcyan},     % 註解顏色
}


% ─────────────────────────────────────────
%  六、數學設定
% ─────────────────────────────────────────
% amsmath：核心數學環境（equation, align, gather 等）
% mathtools：amsmath 的擴充，提供更多對齊與標記工具（已於第四區段載入）
% amssymb：額外數學符號集（\mathbb, \mathfrak, \varnothing 等）
% empheq：強調公式，可為 align 等環境加上框線或括號
% amsthm：定理、引理、命題、定義等環境（\newtheorem）
% cancel：在公式中畫刪除線（\cancel, \bcancel, \xcancel）
% mleftright：修正 \left / \right 在某些環境下產生多餘間距的問題
\usepackage{amsmath, amssymb, empheq, amsthm, cancel, mleftright}

% 自訂數學運算子（顯示為直立羅馬體）
\DeclareMathOperator{\erf}{erf}           % 誤差函數
\DeclareMathOperator{\arcsec}{arcsec}     % 反正割
\DeclareMathOperator{\arccot}{arccot}     % 反餘切
\DeclareMathOperator{\arccsc}{arccsc}     % 反餘割
\DeclareMathOperator{\supp}{supp}         % 支撐集
\DeclareMathOperator{\diag}{diag}         % 對角矩陣
\DeclareMathOperator*{\argmax}{arg\,max}  % argmax（下標置於正下方）
\DeclareMathOperator*{\argmin}{arg\,min}  % argmin（下標置於正下方）

\renewcommand\qedsymbol{$\blacksquare$}   % 將 QED 符號改為實心黑方塊 ■
\allowdisplaybreaks                        % 允許 align 等環境的公式跨頁斷行


% ─────────────────────────────────────────
%  七、超連結設定
% ─────────────────────────────────────────
% hyperref：為目錄、\ref、\cite、URL 等加上可點擊的超連結
%（此套件建議盡量放在所有套件的後段載入以避免衝突）
\usepackage{hyperref}
\hypersetup{
    colorlinks = true,
    linkcolor  = blue,    % 內部連結（目錄項目、\ref 等）
    filecolor  = magenta, % 本機檔案連結
    urlcolor   = blue,    % 外部網址
}
\urlstyle{same}           % URL 使用與正文相同的字型（而非等寬字）


% ─────────────────────────────────────────
%  八、標題樣式設定
% ─────────────────────────────────────────
% titlesec：自訂 \section、\subsection、\subsubsection 的字型大小與格式
\usepackage{titlesec}
\titleformat{\section}      {\Large\CJKfamily{yuan}}     {\thesection}{1em}{}
\titleformat{\subsection}   {\large\CJKfamily{yuan}}     {\thesubsection}{1em}{}
\titleformat{\subsubsection}{\normalsize\CJKfamily{yuan}}{\arabic{section}.\arabic{subsection}.\arabic{subsubsection}\quad}{0em}{}


% ─────────────────────────────────────────
%  九、字型設定（XeLaTeX 專用）
% ─────────────────────────────────────────
% 字型、數學字族與 \kai \yuan \hei 指令由 main.tex 載入的 kelong-cjk.sty 提供（需以 XeLaTeX 編譯）



% ─────────────────────────────────────────
%  十、顏色定義
% ─────────────────────────────────────────
\definecolor{lightgray}    {RGB}{242, 242, 242}   % 淺灰，用於程式碼背景
\definecolor{shadecolor}   {rgb}{0.94, 0.97, 1.0} % 淡藍，用於 framed 陰影底色
\definecolor{sinopia}      {rgb}{0.8,  0.25, 0.04}% 深橙紅
\definecolor{brightmaroon} {rgb}{0.76, 0.13, 0.28}% 亮栗紅
\definecolor{deepcarmine}  {rgb}{0.66, 0.13, 0.24}% 深胭脂紅
\definecolor{darkcyan}     {rgb}{0.0,  0.55, 0.55}% 深青，用於程式碼字串與註解


% ─────────────────────────────────────────
%  十一、頁首頁尾設定
% ─────────────────────────────────────────
% \mytitle 與 \myauthor 由各文件的「0. 標題與作者.tex」以 \def 定義
\pagestyle{fancy}
\fancyhf{}
\lhead{\normalsize\CJKfamily{yuan}\shortstack[l]{\ifx\mychapter\empty\else\mychapter\\\fi\mytitle}}
\rhead{\normalsize\CJKfamily{yuan}\shortstack[r]{\ifx\myinfo\empty\else\myinfo\\\fi\myauthor}}
\cfoot{\footnotesize \thepage} % 頁尾置中：頁碼

% 僅將 subsubsection 標題記錄至頁眉標記（\rightmark），
% section / subsection 不顯示，避免頁首過長
\makeatletter
\renewcommand{\subsubsectionmark}[1]{\markright{#1}}
\renewcommand{\sectionmark}[1]{}
\renewcommand{\subsectionmark}[1]{}
\makeatother


% ─────────────────────────────────────────
%  十二、自訂環境
% ─────────────────────────────────────────
% tcolorbox：高度可自訂的彩色方框套件，以下所有環境皆以此為基礎
\usepackage{tcolorbox}
\tcbuselibrary{skins}   % 載入 skins 函式庫，支援 enhanced 樣式與懸浮標題

% ── whitebox ──────────────────────────────────────────────────
% 無邊框、白色背景的空白方框，作為 narrowframed 的底層元件，
% 也可直接使用以產生視覺隔離但不顯眼的區塊。
%
% 用法：
%   \begin{whitebox}
%     內容...
%   \end{whitebox}
% ──────────────────────────────────────────────────────────────
\newtcolorbox{whitebox}{colframe=white}

% ── narrowframed ──────────────────────────────────────────────
% 縮排至頁寬 90%、水平置中的方框，底色為 shadecolor（淡藍），
% 適合放置定義、定理、重要結論等需要視覺凸顯的內容。
%
% 用法：
%   \begin{narrowframed}
%     \textbf{定義 X.X} 定義名稱\\
%     定義內容...
%   \end{narrowframed}
% ──────────────────────────────────────────────────────────────
\newenvironment{narrowframed}
    {\begin{center}\begin{minipage}{0.9\textwidth}\begin{whitebox}}
    {\end{whitebox}\end{minipage}\end{center}}

% ── examplebox ────────────────────────────────────────────────
% 黑框題目方塊，左上角帶有灰色懸浮標籤（顯示題號與名稱），
% 適合放置習題、例題或小考題目。
%
% 參數：
%   #1 ： 題號，顯示為「習題 #1」
%   #2 ： 題目名稱
%
% 用法：
%   \begin{examplebox}{1.1}{題目名稱}
%     題目敘述...
%   \end{examplebox}
% ──────────────────────────────────────────────────────────────
\newtcolorbox{examplebox}[2]{
    enhanced,
    colback  = gray!5,
    colframe = black,
    boxrule  = 1.5pt,
    arc      = 3mm,
    width    = 0.9\textwidth,
    center,
    left=5mm, right=5mm, top=5mm, bottom=4.5mm,
    attach boxed title to top left={yshift=-3mm, xshift=5mm},
    boxed title style={
        enhanced,
        colback  = gray!60,
        colframe = black,
        boxrule  = 1.5pt,
        arc      = 2mm,
        left=3mm, right=3mm, top=1mm, bottom=1mm,
    },
    fonttitle = \bfseries,
    title     = {習題 #1\;\; #2}
}

% ── solutionbox ───────────────────────────────────────────────
% 解答環境，無框、無背景，前方自動產生「解答」粗體標題。
% 通常緊接在 examplebox 之後使用。
%
% 用法：
%   \begin{solutionbox}
%     解題過程...
%   \end{solutionbox}
% ──────────────────────────────────────────────────────────────
\newenvironment{solutionbox}{%
    \par\vspace{0.5em}%
    {\normalsize\textbf{解答}}\\[0.4em]%
}{%
    \par\vspace{1em}%
}

% ── analysisbox ───────────────────────────────────────────────
% 題目分析環境，以 narrowframed 相同的淡藍底框呈現，
% 框內頂端自動顯示「題目分析」粗體標題。
% 通常接在 solutionbox 之後，補充解題思路、考點或觀念說明。
%
% 用法：
%   \begin{solutionbox}
%     解題過程...
%   \end{solutionbox}
%   \begin{analysisbox}
%     這題考的是... 解題關鍵在於...
%   \end{analysisbox}
% ──────────────────────────────────────────────────────────────
\newenvironment{analysisbox}{%
    \par\vspace{0.5em}%
    \begin{center}\begin{minipage}{0.9\textwidth}\begin{whitebox}%
    \noindent{\normalsize\textbf{題目分析}}\\[0.4em]%
}{%
    \end{whitebox}\end{minipage}\end{center}%
    \par\vspace{1em}%
}


% ─────────────────────────────────────────
%  十三、雜項設定
% ─────────────────────────────────────────
% gensymb：提供通用符號，如 \degree（°）、\celsius（℃）、\micro（µ）
\usepackage{gensymb}

% parskip：以段落間距（垂直空白）取代首行縮排
\usepackage{parskip}
\parindent = 0pt

% setspace：調整行距；\onehalfspacing 設為 1.5 倍行距
\usepackage{setspace}
\onehalfspacing

% enumitem：自訂 itemize / enumerate 清單的符號、縮排與間距
\usepackage{enumitem}

% ragged2e：提供 \RaggedRight 等改進版對齊指令，斷行比原生版本更均勻
\usepackage{ragged2e}

% changepage：用於局部臨時調整版面寬度（如縮排區塊）
\usepackage{changepage}

% multicol：多欄排版，\begin{multicols}{2}...\end{multicols}
\usepackage{multicol}
